\setcounter{deso}{1}
\begin{name}
	{\tenchude}
	{\tendethi}
	{\tentruong}
	{\thoigian}
\end{name}
\begin{bt}[1,5 điểm]
\begin{enumerate}[1.]
    \item Tính giá trị biểu thức $A=3\sqrt{2}-\sqrt{8}+2\sqrt{18}-5\sqrt{50}$.
    \item Giải phương trình $16x^2-1=0$.
    \item Giải bất phương trình $x<2-2x$.
\end{enumerate}
\end{bt}

\begin{bt}[2,0 điểm]
\begin{enumerate}[1.]
    \item Cho hàm số $y=ax^2$ với $a \ne 0$ có đồ thị đi qua điểm $M(4;2)$. Xác định hệ số $a$.
    \item Cho biểu thức $M=\dfrac{\sqrt{x}}{\sqrt{x}+1}+\dfrac{2}{\sqrt{x}-1}-\dfrac{x+3\sqrt{x}}{x-1}$
    \begin{enumerate}
        \item Rút gọn biểu thức $M$.
        \item Tìm tất cả giá trị nguyên của $x$ để $M$ nhận giá trị nguyên.
    \end{enumerate}
\end{enumerate}
\end{bt}

\begin{bt}[1,5 điểm]%[Đề tham khảo tuyển sinh vào 10 - Quận Gò Vấp - Đề số 1]%[9D2T3]%[TS10B7]
Một chiếc xe buýt chạy tuyến bến xe An Sương - Suối Tiên - Đại học Quốc Gia thành phố Hồ Chí
Minh có sức chứa tối đa $80$ chỗ. Trong xe hiện có $75 \%$ số hành khách đang đứng và $\dfrac{1}{9}$ số hành khách là các cụ già. Hỏi xe buýt này đang chở bao nhiêu hành khách? Biết rằng có $38$ sinh viên của đại học Quốc Gia thành phố Hồ Chí Minh cũng đang ở trong xe.
\loigiai{
Gọi $x$ là số hành khách trên xe. \\
Số hành khách đứng (gồm sinh viên và người khác): $0,75x$\\
Số hành khách ngồi (gồm cụ già và sinh viên): $0,25x$\\
Số cụ già ngồi: $\dfrac{1}{9}x$\\
Số sinh viên ngồi: $0,25x-\dfrac{1}{9}x=\dfrac{5}{36}x$\\
Số sinh viên đứng: $38-\dfrac{5}{36}x$\\
Số người đứng còn lại: $y$\\
Ta có phương trình:$38-\dfrac{5}{36}x+y=0,75$ $\Leftrightarrow y=\dfrac{8}{9}x-38$\\
Mà: $38 < x\le 80$ và $x\vdots 9,x\vdots 36$ $\Rightarrow x=72$ .\\
Vậy số hành khách trên xe buýt là $72$ .
}
\end{bt}

\begin{bt}[1 điểm]%[Đề tham khảo tuyển sinh vào 10 - Quận Gò Vấp - Đề số 1]%[8H4T2]%[TS10B6]
Một lọ vitamin dạng hình trụ với bán kính đáy là ${R}=1,5 {cm}$ và chiều cao là ${h}=10 {~cm}$, nắp lọ được thiết kế đặc biệt gắn liền với một ống hình trụ cao $2$ cm để chứa được các hạt hút ẩm (Silicagel), khi đóng nắp nó sẽ giúp bảo vệ các viên sủi vitamin khỏi ẩm mốc và chảy nước. Những viên sủi được xếp đầy trong lọ cũng có dạng hình trụ với diện tích đáy bằng với diện tích đáy lọ và thể tích của mỗi viên là $1,8 \pi\left(\mathrm{cm}^{3}\right)$.
\begin{listEX}
\item Hỏi trong lọ chứa bao nhiêu viên sủi vitamin? Biết công thức tính thể tích hình trụ là $V=\pi R^{2} h$.
\item Những lọ vitamin này được xếp thẳng đứng, sát nhau vào một hộp bao bì có dạng hình hộp chữ nhật. Hỏi thể tích của hộp bao bì là bao nhiêu để có thể chứa được $20$ lọ xếp thành $5$ hàng, mỗi hàng $4$ lọ ? ( bỏ qua bề dày của giấy bao bì )
\end{listEX}
\dapso{$V_{hop}=15.12.10=1800\left(c{m^3}\right)$}
\loigiai{~
\begin{enumerate}
\item Chiều cao mỗi viên sủi vitamin: $h=\dfrac{V}{\pi{R^2}}=\dfrac{1,8\pi}{\pi .1,5^2}=0,8\left(cm\right)$\\
Số viên sủi vitamin mà lọ có thể chứa được: $\dfrac{10-2}{0.8}=10$ (viên)
\item Chiều dài của hộp bao bì: $5.1,5.2=15\left(cm\right)$\\
Chiều rộng của hộp bao bì: $4.1,5.2=12\left(cm\right)$\\
Thể tích của hộp bao bì: $V_{hop}=15.12.10=1800\left(c{m^3}\right)$
\end{enumerate}
}
\end{bt}

\begin{bt}[1,0 điểm]
Một hộp đựng 25 viên bi đỏ được đánh số từ 1 đến 25 và 26 viên bi xanh được đánh số từ 1 đến 26. Xét phép thử "Chọn ngẫu nhiên 1 viên bi trong hộp". Viết không gian mẫu của phép thử và tính xác suất của biến cố $A$: "Chọn được viên bi đánh số chia hết cho 5.
\end{bt}

\begin{bt}[3,0 điểm]%[Đề tham khảo tuyển sinh vào 10 - Bình Chánh - Đề số 2]%[9H3G9]%[TS10K10]
Từ điểm $ A $ ngoài đường tròn $ (O; R) $ vẽ hai tiếp tuyến $ AB $, $ AC $ với đường tròn $ (O) $ ($ B $, $ C $ là tiếp điểm). Gọi $ H $ là giao điểm của $ AO $ và $ BC $. Gọi $ I $ trung điểm của $ AB $. Từ $ B $ kẻ đường thẳng vuông góc với $ OI $ tại $ K $, đường thẳng này cắt đường tròn $ (O) $ tại $ D $ ($ D $ khác $ B $).
\begin{listEX}
\item Chứng minh: Tứ giác $ ABOC $ nội tiếp và $ OK\cdot OI = OH\cdot OA $.
\item Đường tròn $\left(I;\dfrac{AB}{2}\right)$ cắt $ AC $ tại $ E $. Gọi $ F $ là giao điểm của $ BE $ và $ OA $. Chứng minh: $ F $ đối xứng với $ O $ qua $ H $?
\item Chứng minh rằng: Đường tròn ngoại tiếp $\triangle AFB$ đi qua điểm $ K $?
\end{listEX}
\loigiai{
\begin{center}
\begin{tikzpicture}[scale=0.7,>=stealth, font=\footnotesize, line join=round, line cap=round]
\path
(0,0) coordinate (O)
(120:3) coordinate (B)
(-120:3) coordinate (C)
;
\tkzDefLine[orthogonal=through B](O,B)\tkzGetPoint{d}
\tkzDefLine[orthogonal=through C](O,C)\tkzGetPoint{e}
\tkzInterLL(B,d)(C,e)\tkzGetPoint{A}
\coordinate (I) at ($(A)!.5!(B)$);
\draw[red] (O) circle(3cm);
\tkzDrawCircle[radius,blue](I,A)
\tkzInterLL(A,O)(B,C)\tkzGetPoint{H}
\tkzDefLine[orthogonal=through B](O,I)\tkzGetPoint{f}
\tkzInterLL(B,f)(I,O)\tkzGetPoint{K}
\tkzInterLC[R](B,f)(O,3cm)\tkzGetPoints{D}{D'}
\tkzInterLC(A,C)(I,A)\tkzGetPoints{E}{E'}
\tkzInterLL(A,O)(B,E)\tkzGetPoint{F}
\draw
(A)--(B)--(C)--(O)--cycle
(I)--(O)
(C)--(O)--(B)--(D)
(A)--(C)
(B)--(E) (A)--(K)--(H)
;
\foreach \x/\y/\z in {O/C/A,A/B/O,O/H/B,I/K/D,B/E/A}\draw pic[draw,angle radius =2mm] {right angle = \x--\y--\z};
\foreach \x/\g in
{A/180,B/90,C/-90,O/0,I/90,H/-60,K/30,D/-20,E/-90,F/-150}
\fill[black](\x) circle (2pt)
($(\x)+(\g:5mm)$) node{$\x$};
\end{tikzpicture}
\end{center}
\begin{enumEX}{1}
\item Xét tứ giác $ABOC$ có $\widehat{ABO}=90{}^\circ $ ($AB$ là tiếp tuyến);
$\widehat{ACO}=90{}^\circ $ ($AC$ là tiếp tuyến).\\
$\Rightarrow \widehat{ABO}+\widehat{ACO}=180{}^\circ $.\\
$\Rightarrow $Tứ giác $ABOC$ nội tiếp (tổng hai góc đối bằng $180{}^\circ $) .\\
\textit{Chứng minh: $OK\cdot OI\text{ }=\text{ }OH\cdot OA$.}\\
Xét $\triangle IBO$ vuông tại $B$ có
\begin{itemize}
\item Đường cao $BK \Rightarrow OB^2=OK\cdot OI$.
\item $AB=AC$ (tính chất 2 tiếp tuyến cắt nhau).
\item  $OB=OC$ (bán kính).
\end{itemize}
$\Rightarrow $ $OA$ là đường trung trực của $BC$.\\
Xét $\triangle IBO$ vuông tại $B$ có đường cao $BH \Rightarrow OB^2=OH\cdot OA$
 $\Rightarrow OK\cdot OI=OH\cdot OA \left( =OB^2\right)$.
\item
Ta có $\widehat{AEB}=90{}^\circ $ (góc nội tiếp chắn nửa đường tròn) $\Rightarrow \widehat{FEC}=90{}^\circ $ (góc kề bù).\\
$\Rightarrow \widehat{FHC}=90{}^\circ $ ($AO\perp BC$) $\Rightarrow \widehat{FEC}+\widehat{FHC}=180{}^\circ $, suy ra $CEFH$ nội tiếp.\\
$\Rightarrow \widehat{BCA}=\widehat{BFO}$ (góc ngoài và góc đối trong).\\
Mà $\widehat{BOA}=\dfrac{1}{2}\widehat{BOC}=\dfrac{1}{2}$ sđ $\wideparen{BC}=\widehat{BCA}$.\\
$\widehat{BFO}=\widehat{BOF}=\left( \widehat{BCA} \right)\Rightarrow \triangle BFO$ cân tại $B$. \\
Có $BH$ là đường cao nên là đường trung trực của $FO$.\\
Vậy $F$ đối xứng với $O$ qua điểm $H$.
\item Ta có $OK\cdot OI=OH\cdot OA$ $\Rightarrow \dfrac{OK}{OA}=\dfrac{OH}{OI}$, $\widehat{AOI}$ chung\\ $\Rightarrow $ $\triangle KOH\backsim \triangle AOI$ (c.g.c) $\widehat{OAI}=\widehat{OKH}\Rightarrow $ tứ giác $AIKH$ nội tiếp (góc ngoài = góc đối trong).\\
$\Rightarrow \widehat{IKA}=\widehat{IHA}$ ($2$ góc nội tiếp cùng chắn cung $IA$). Mà:
\begin{itemize}
\item $\widehat{IHA}=\widehat{IAH}$ ($\triangle IAH$ cân tại $I$).
\item $\widehat{IAH}=\widehat{OBH}$ (cùng phụ $\widehat{BOA}$). \item $\widehat{OBH}=\widehat{FBH}$ ($BH$ là phân giác).
\end{itemize}
 $\Rightarrow \widehat{IKA}=\widehat{FBH}$.\\
Mặt khác $\widehat{BKI}=\widehat{BFH}=90{}^\circ \Rightarrow \widehat{IKA}+\widehat{BKI}=\widehat{FBH}+\widehat{BHF}\Leftrightarrow \widehat{BKA}=\widehat{BFA}$ (góc ngoài $\triangle BHF$).\\
$\Rightarrow $ Tứ giác $ABKF$ nội tiếp ($2$ đỉnh liên tiếp cùng nhìn $1$ cạnh).\\
$\Rightarrow $ Đường tròn ngoại tiếp $\triangle ABF$ đi qua điểm $K$.
\end{enumEX}
}
\end{bt}